\documentclass[11pt]{article}
\usepackage[margin=1in]{geometry}
\usepackage{amsmath,amssymb,booktabs,xcolor,tcolorbox}
\newcommand{\1}{\mathbf{1}}
\newtcolorbox{idea}{colback=blue!4,colframe=blue!40,title=In plain words}
\title{Simplified proof 1b: the covariance law on one concrete sequence}
\author{Folder 22 / simplified\_proofs --- companion to \texttt{1\_covariance\_law.pdf}}
\date{}
\begin{document}\maketitle

\section{The sequence}
Real tokens are 0--3, noise tokens are 4--7. The model stands at position $i=12$ and must predict
position 13. The true next token is $y=3$.
\[
\begin{array}{c|ccccccccccc|c}
j & 1&2&3&4&5&6&7&8&9&10&11 & 12\ (\text{query})\\\hline
x_j & 2&3&1&2&0&6&2&5&3&1&2 & 4
\end{array}
\]

\section{The three ingredients for every position $j$}
\begin{itemize}
\item $\1[x_j = y]$ (the \textbf{indicator function}) $=1$ if $x_j = y$ (the token there is 3), else $0$.
\item $\varphi_j = \dfrac{\partial s_j}{\partial \kappa}$: what the chosen weight $\kappa$ multiplies in position $j$'s score
\[
s_j = a_2\,\nu_j + b_2\,\sigma_i\,r_j[x_i] + c_2\,\langle r_i, r_j\rangle .
\]
\item $A_j$: the attention on position $j$. To keep the numbers simple, take \textbf{equal attention}:
$A_j = \tfrac{1}{11}$ for every $j$ (this is what the model does at the start of training).
\end{itemize}

Here $r_j$ is ``the last real token before $j$'' (layer 1, once it has learned), and $r_{12}=2$.
\[
\begin{array}{l|ccccccccccc}
j & 1&2&3&4&5&6&7&8&9&10&11\\\hline
x_j & 2&3&1&2&0&6&2&5&3&1&2\\
r_j & -&2&3&1&2&0&0&2&2&3&1\\\hline
\1[x_j=y] & 0&\mathbf 1&0&0&0&0&0&0&\mathbf 1&0&0\\
\varphi_j\ \text{for } a_2\ (=\nu_j) & 0&0&0&0&0&\mathbf 1&0&\mathbf 1&0&0&0\\
\varphi_j\ \text{for } b_2\ (=\sigma_{12}\,r_j[4]) & 0&0&0&0&0&0&0&0&0&0&0\\
\varphi_j\ \text{for } c_2\ (=\langle r_{12},r_j\rangle) & 0&\mathbf 1&0&0&\mathbf 1&0&0&\mathbf 1&\mathbf 1&0&0
\end{array}
\]
($\varphi_j=0$ for $b_2$ at every $j$ because the query token 4 is noise, so $\sigma_{12}=0$.)

\section{The formula we start from}
Steps 1--6 of \texttt{1\_covariance\_law.pdf} give
\[
-\frac{\partial \ell}{\partial \kappa} \;=\; w \sum_{j} A_j\,\big(\1[x_j=y] - f[y]\big)\,\varphi_j,
\qquad
f[y] = \sum_j A_j\,\1[x_j=y],
\qquad
w = \frac{1}{f[y] + b_y}.
\]
Here: $f[y] = 2\cdot\tfrac1{11} = \tfrac{2}{11}$ and $w = 1/(\tfrac2{11}+0.01)\approx 5.2$.

\section{What ``average'' and ``covariance'' mean here}
Pick one position $J$ at random, with probability $A_J$. For any quantity $g$,
\[
\mathbb E[g] \;=\; \sum_j A_j\, g_j \qquad\text{(the attention-weighted average).}
\]
So $f[y] = \mathbb E[\1[x_J=y]]$. The covariance of two quantities is
\[
\mathrm{Cov}(\1[x_J=y],\varphi) \;=\; \mathbb E[\1[x_J=y]\cdot\varphi] \;-\; \mathbb E[\1[x_J=y]]\;\mathbb E[\varphi].
\]
\begin{idea}
$\mathbb E[\1[x_J=y]\cdot\varphi]$ = how often a position is \emph{both} right \emph{and} pushed up by $\kappa$.\\
$\mathbb E[\1[x_J=y]]\,\mathbb E[\varphi]$ = how often that would happen \emph{by pure chance}.\\
Covariance $>0$: the weight pushes up right positions more than chance. $<0$: less than chance.
\end{idea}

\section{Why the gradient sum IS the covariance}
Write the sum as an average:
\[
\sum_j A_j\,(\1[x_j=y] - f[y])\,\varphi_j
= \sum_j A_j\,\1[x_j=y]\,\varphi_j \;-\; f[y]\sum_j A_j\,\varphi_j
= \mathbb E[\1[x_J=y]\cdot\varphi] - \mathbb E[\1[x_J=y]]\,\mathbb E[\varphi]
= \mathrm{Cov}(\1[x_J=y],\varphi).
\]
(First equality: multiply out the bracket. Second: $\sum_jA_j(\cdot)=\mathbb E[\cdot]$ and $f[y]=\mathbb E[\1[x_J=y]]$.) Hence
\[
\boxed{\;-\frac{\partial \ell}{\partial \kappa} \;=\; w\cdot \mathrm{Cov}(\1[x_J=y],\varphi)\;}
\]

\section{The numbers, weight by weight}
\renewcommand{\arraystretch}{1.6}
\[
\begin{array}{l|c|c|c|c|c}
\kappa & \mathbb E[\1[x_J=y]] & \mathbb E[\varphi] & \mathbb E[\1[x_J=y]\cdot\varphi] & \mathrm{Cov} & \text{GD moves }\kappa\\\hline
c_2 & \tfrac{2}{11} & \tfrac{4}{11}\ (\text{pos }2,5,8,9) & \tfrac{2}{11}\ (\text{pos }2,9) & \tfrac{2}{11}-\tfrac{8}{121}=+\tfrac{14}{121} & \text{up}\\
a_2 & \tfrac{2}{11} & \tfrac{2}{11}\ (\text{pos }6,8) & 0 & 0-\tfrac{4}{121}=-\tfrac{4}{121} & \text{down}\\
b_2 & \tfrac{2}{11} & 0 & 0 & 0 & \text{no move}
\end{array}
\]

\paragraph{Check $c_2$ with the sum directly.} Only $j$ with $\varphi_j=1$ contribute:
\[
\tfrac1{11}\Big[\underbrace{(1-\tfrac2{11})}_{j=2}+\underbrace{(0-\tfrac2{11})}_{j=5}+\underbrace{(0-\tfrac2{11})}_{j=8}+\underbrace{(1-\tfrac2{11})}_{j=9}\Big]
=\tfrac1{11}\cdot\tfrac{9-2-2+9}{11}=\tfrac{14}{121}. \checkmark
\]
\paragraph{Check $a_2$.}
\[
\tfrac1{11}\Big[\underbrace{(0-\tfrac2{11})}_{j=6}+\underbrace{(0-\tfrac2{11})}_{j=8}\Big]=-\tfrac{4}{121}. \checkmark
\]
So on this query: $-\partial\ell/\partial c_2 = 5.2\times\tfrac{14}{121}\approx +0.60$ (raise $c_2$), and
$-\partial\ell/\partial a_2 = 5.2\times(-\tfrac{4}{121})\approx -0.17$ (lower $a_2$, i.e.\ look at noise less).

\begin{idea}
The covariance law is one formula used three times. Each weight has its own $\varphi$ (what it multiplies in
the score). Gradient descent raises a weight exactly when the positions it pushes up hold the right answer
more often than chance.
\end{idea}
\end{document}
